\section{Introduction}\label{sec:intro}
\paragraph{Background.} Road traffic injuries remain a major public-health problem: the World Health Organization estimates 1.19 million road-traffic deaths per year, of which pedestrians account for 274,000 deaths (23\%)~\citep{who2023road}. Automated driving systems and advanced driver-assistance systems (ADAS) must therefore interact safely with pedestrians, avoiding collisions and allowing crossings to take place without hesitation~\citep{iddped2025}. A central capability is the anticipation of pedestrian motion: a vehicle that can estimate where a pedestrian will be one to two seconds ahead can brake, steer or yield early, whereas a purely reactive system must respond after the pedestrian has already entered the vehicle's path. Pedestrian trajectory prediction, i.e.\ forecasting the future bounding boxes of a pedestrian from an egocentric (vehicle-mounted) view, is accordingly a core component of the perception-and-planning stack.

\paragraph{Problem context.} Most learning-based predictors have been developed and evaluated on datasets collected in structured traffic, such as PIE and JAAD, in which crosswalks, traffic signals and yielding behaviour regularise pedestrian motion~\citep{rasouli2019pie,iddped2025}. These conditions are not representative of many parts of the world. In unstructured traffic, pedestrians walk along carriageways, cross diagonally between moving vehicles, hesitate or change direction after beginning to cross, and share the road with a dense and heterogeneous mix of vehicles, often without signals or marked crossings~\citep{iddped2025}. The recently released IDD-PeD dataset captures such conditions in Indian traffic and reports that established trajectory predictors degrade markedly on it: the error of the Multimodal Transformer Network (MTN) increases by up to 1,208 mean-squared-error units relative to its error on PIE~\citep{iddped2025}. Reliable prediction in unstructured traffic therefore remains an open problem.

\paragraph{Challenges.} Three difficulties are especially relevant. First, the \emph{motion cues are weak and ambiguous}: when pedestrians hesitate, pause or change direction, the past trajectory of a bounding box carries little information about the intended path, so predictors that rely on box history, ego-vehicle speed and optical flow alone have limited evidence on which to base a forecast. Second, \emph{pedestrians are frequently occluded and the scene varies strongly}, which further degrades visual and motion evidence and makes the appearance of the scene a poor guide to behaviour. Third, \emph{scene knowledge exists but is not exploited}: datasets such as IDD-PeD annotate, for every pedestrian track, the intersection type, signalisation, road type, location type, motion direction and occlusion state~\citep{iddped2025}. These attributes distinguish, for example, a pedestrian waiting at a signalised junction from one walking along an undivided road, yet compact trajectory models such as MTN~\citep{yin2021mtn} do not use them. Methods that do combine behavioural cues with scene information typically rely on image-level features such as regions of interest or segmentation maps and on stochastic decoders, and are evaluated on structured-traffic benchmarks~\citep{rasouli2022pedformer,sharma2025dual}.

\paragraph{Motivation.} The availability of cheap, interpretable, native annotations motivates a lightweight use of scene context. Two practical questions arise. (i)~\emph{How should static context be injected into an attention model without disturbing the motion pathway that has already been learned, and without harming samples for which the context is uninformative?} A single context token that a gate may admit or reject for each sample offers a safe fallback to the original model. (ii)~\emph{Can the additional capacity be paid for by removing redundancy elsewhere in the network?} Compact attention-based predictors are small models trained on limited data, so additional modules increase the risk of over-fitting, while their stacked encoder layers repeat the same type of self-attention computation. Sharing parameters across layers reduces the parameter count and can act as a regulariser~\citep{dehghani2019universal,lan2020albert}, but tying entire layers is harmful when layers play different roles, so the question becomes which sub-modules can be tied safely. Finally, because the study concerns small differences between closely related models, a controlled comparison under one protocol, with several seeds, is needed to separate genuine effects from run-to-run variation.

\paragraph{Approach and contributions.} This paper proposes PACE-Net, short for \emph{Pedestrian trajectory prediction with Admission-gated Context Enrichment} (Fig.~\ref{fig:overall}), a new architecture for pedestrian trajectory prediction in unstructured traffic. Its design rests on a separation of concerns: a \emph{motion pathway} models how the pedestrian moves, a \emph{context pathway} describes where the pedestrian is, and the two are merged only at the decoder memory, under the control of a learned admission gate. PACE-Net is built from four named components, which constitute the contributions of this paper.
\begin{itemize}[leftmargin=1.4em,itemsep=1pt,topsep=2pt]
\item \textbf{ROUTE (Routed Operator Unification of Trajectory and Environment).} An adaptive context-memory fusion module in which a selector chooses, for each sample, among three ways of merging scene context with the trajectory memory (identity, additive broadcast and cross-attention), so that the fusion strategy itself is learned (Sec.~\ref{ssec:fusion}).
\item \textbf{SAGE (Sample-wise Admission Gate for Evidence).} A hard, per-sample gate trained with a straight-through Gumbel-softmax decides whether the scene context is used at all, with the motion-only prediction as an exact fallback (Sec.~\ref{ssec:gate}).
\item \textbf{MASC (Memory Augmentation with Source-masked Context).} The gated 32-D native scene-context token, built from five IDD-PeD scene labels and the observed occlusion history, is appended to the decoder memory ($15\!\to\!16$ tokens) and masked out of attention when rejected (Sec.~\ref{ssec:memory}).
\item \textbf{SATIE (Self-Attention Tying in the Encoder).} A selective weight-parameter sharing scheme in which only the self-attention projections $(W_Q,W_K,W_V,W_O)$ of the three encoder layers are tied, reducing the parameter count by $6.07\%$ while all other modules remain layer-specific (Sec.~\ref{ssec:c4}).
\end{itemize}
These components are embedded in a dual-pathway architecture with its own design rationale (Sec.~\ref{ssec:rationale}) and are evaluated in a controlled, matched study, including a $2{\times}2$ factorial comparison of the two context-integration strategies with and without sharing, with per-seed numbers and paired statistics. In that study, ROUTE did not improve over the SAGE-gated token, and the final configuration therefore combines SAGE, MASC and SATIE; the negative result is reported openly.
The experiments constitute a proof-of-concept study (24 test trajectories, one training epoch, Farneb\"ack optical flow). The best model, \CAMTNS, has the lowest mean error and the fewest parameters, but its advantage over \CAMTN is within seed variance. The study is therefore reported as evidence for further validation and not as a benchmark result. The remainder of the paper is organised as follows: Section~\ref{sec:related} reviews related work, Section~\ref{sec:method} presents the method, Section~\ref{sec:results} the experiments, Section~\ref{sec:ablation} the ablation study, and Section~\ref{sec:limits} the limitations and conclusions.
